\chapter{家庭、服务、医疗与康复}
\label{chap:service-healthcare-case}

\section{辅助进食是完整服务系统，不是抓取演示}

辅助进食同时涉及食物感知、叉取/舀取、无碰转移、入口定位、用户时机、卫生和社会互动。失败可能是没有获取食物、食物掉落、用户咳嗽、头部移动、餐盘移位或交互设备不可用。与仓内抓取不同，最终执行对象靠近人的面部，速度、力、工具形状和用户控制权都是核心需求。

一次入口成功是多个阶段串联：
\begin{equation}
P_{\mathrm{bite}}=P(A)P(T\mid A)P(D\mid A,T)P(U\mid A,T,D),
\label{eq:feeding-chain-success}
\end{equation}
其中 $A$ 是成功获取食物，$T$ 是安全转移，$D$ 是正确递送到口部，$U$ 是用户在合适时机接受。即使每阶段成功率均为 0.95，串联成功率也只有约 0.81；因此必须设计中途检测和恢复，而不能只优化叉取。

\section{从 PR2 实验到院外系统}

Park 等的主动辅助进食系统使用 PR2、视觉引导行为、用户界面和安全工具，并由 10 名无运动障碍参与者和 9 名运动障碍参与者评价\cite{park2020feeding}。该研究把“机器人把食物送到嘴边”展开为用户选择、视觉定位、主动靠近和可感知控制，是比单组件精度更强的端到端用户证据。

后续院外系统进一步改变硬件和自主组织。Nanavati 等与运动障碍社区研究者共同设计便携系统，在办公室、食堂等地点由 5 名用户和一名社区研究者完成用餐，并由一名社区研究者在家中连续 5 天使用 10 餐；论文报告总计超过 15 小时的真实环境使用\cite{nanavati2025feeding}。其关键经验不是“更强模型消除异常”，而是可定制界面、用户参与选择和可变自主帮助处理异常。

\begin{figure}[H]
\centering
\begin{tabular}{c@{\ $\rightarrow$\ }c@{\ $\rightarrow$\ }c@{\ $\rightarrow$\ }c}
\fbox{用户选食物/方式} & \fbox{获取与确认} & \fbox{转移/入口时机} & \fbox{接受、取消或恢复}
\end{tabular}
\caption{辅助进食状态链。用户始终保留选择、暂停和取消能力。}
\label{fig:feeding-state-chain}
\end{figure}

\begin{table}[H]
\begingroup\hbadness=10000
\centering
\caption{证据从组件和实验室走向院外，但尚未等于产品交付}
\label{tab:feeding-evidence-ladder}
\begin{tabularx}{\textwidth}{p{0.18\textwidth}YYYp{0.20\textwidth}}
\toprule
层级 & 主要问题 & 代表证据 & 新暴露的失效 & 仍缺证据 \\
\midrule
组件试验 & 食物检测/获取是否可行 & 单食品、重复 trial & 类别与形变长尾 & 整餐、用户和清洁 \\
实验室用户研究 & 人能否控制并接受系统 & PR2、两类参与者 & 头部、界面、舒适性 & 家庭空间与长期使用 \\
院外/家庭研究 & 异常和情境能否处理 & 多地点与 5 天部署 & 座姿、活动、社交、布置 & 大样本、月级可靠性 \\
可持续产品 & 能否安全持续服务 & 尚需认证与运营材料 & 维护、责任、支付、培训 & 事故率、成本与服务网络 \\
\bottomrule
\end{tabularx}
\endgroup
\end{table}

\section{可变自主不是“不够自动”的妥协}

完全自主要求系统识别所有食物、用户状态和社交时机，开放世界代价很高。共享自主把机器擅长的轨迹、避障和稳定递送与用户擅长的偏好、时机和异常解释组合。用户介入率
\begin{equation}
I_u=\frac{N_{\mathrm{user\ corrections}}}{N_{\mathrm{bite\ attempts}}}
\label{eq:user-intervention-rate}
\end{equation}
不应简单越低越好：主动选择和取消体现控制权；被迫频繁纠错才是负担。评测需把意愿性交互与故障补救分开，并报告任务负荷、独立性和尊严等用户结果。

系统应提供机械限力/可脱离工具、速度限制、口部附近专用轨迹、急停、用户暂停和护理者接管。视觉或模型置信度不能取代这些机制。清洁和食物安全、设备安装、坐姿差异、误吸风险和紧急处置需要专业流程；研究原型的伦理审批也不等同于医疗器械准入。

\section{隐私、责任与成本进入系统定义}

面部、语音、饮食和健康数据高度敏感。最小化采集、端侧处理、访问控制、保留期限和删除机制应在设计期确定。远程支持若能查看视频或控制机械臂，必须记录授权、身份、动作范围和审计日志。

责任按阶段分配：模型提出食物或动作候选，运行时保证约束，用户决定偏好和时机，护理者处理超出能力边界的情况，供应方负责维护和更新。成本不能只算机械臂，还包括安装、清洁耗材、培训、远程支持、故障上门和替代照护。若系统每天只节省少量护理时间，却需要高频技术支持，独立性价值可能存在，商业可持续性仍未证明。

\section{知识小结与案例判断}

辅助进食把感知、操作、用户界面、安全和异常恢复绑定在一条靠近人体的服务链。院外研究比实验室成功率更能暴露空间和社会情境，但样本与时长仍有限。案例判断是：服务机器人最重要的泛化机制往往不是完全自动，而是让用户以低负担控制系统、识别异常并安全恢复；产品结论还需要长期可靠性、准入、责任和单位服务成本证据。
